\documentclass{article}
\usepackage[english]{babel}
\usepackage[a4paper,top=2cm,bottom=2cm,left=2.2cm,right=2.2cm]{geometry}
\usepackage{amsmath,amssymb}
\usepackage[hidelinks]{hyperref}
\usepackage{booktabs}
\usepackage{array}
\sloppy

\title{Photon flux}
\author{}
\date{\today}

\begin{document}
\begingroup\let\newpage\relax   
\makeatletter\renewcommand\@maketitle{\begin{center}{\LARGE\@title\par}\vskip 0.5em{\large\@date\par}\end{center}\vskip 0.5em}\makeatother
\maketitle
\endgroup


The flux is written in 
\texttt{src/photon\_flux.cpp} and given as $dN/d\omega$: number of photons
per unit photon energy.


\section{Variables}

The code works with the light-cone momentum of the photon $q^+$, the photon energy $\omega$ and the energy fraction $z_\gamma$ (the fraction
of the beam energy carried by the photon, the Starlight tables call it $y$).
They are related by
\begin{equation}
  \omega = \frac{q^+}{\sqrt2}, \qquad
  z_\gamma = \frac{2\omega}{\sqrt{s_{NN}}}, \qquad\text{so}\qquad
  \frac{dN}{d\omega} = \frac{2}{\sqrt{s_{NN}}}\,\frac{dN}{dz_\gamma},
\end{equation}
with $\sqrt{s_{NN}} = 5.36$~TeV for Pb+Pb. The flux can also be given as $dN/dz_\gamma$, or $z_\gamma\,dN/dz_\gamma$.

\section{Different fluxes in this code}

The flux is chosen with the environment variable \texttt{FLUX\_MODEL}
(the defaults are in \texttt{run\_scripts/config.sh}):

\begin{itemize}
\item \texttt{EFF} (default): ``effective'' flux of arXiv:2404.09731 (Eq.~4).
\item \texttt{STARLIGHT}: table made by P.~Paakkinen in arXiv:2404.09731 (only AnAn and An0n).
\item \texttt{PL}: flux of a point-like nucleus at a fixed impact parameter $b$.
\item \texttt{WS}: flux for a Woods--Saxon charge distribution.
\end{itemize}

The impact parameter $b$ is the distance between the centres of the two
nuclei. The \texttt{PL} and \texttt{WS} depends on $b$, which is integrated by the code. In \texttt{EFF} and \texttt{STARLIGHT} this
integral is already done, so the flux depends only on the photon energy.

Two more settings change the flux:
\begin{itemize}
\item \texttt{CHANNEL}, the neutron class. The nuclei can break up
  and emit neutrons, which is measured in the zero degree calorimeters. The
  probability that a nucleus does not break up is
  $P_\text{noEM}(b) = \exp(-S/b^2)$ with $S = (17.4~\text{fm})^2$. The flux is
  multiplied by $1$ for AnAn (any number of neutrons), by $P_\text{noEM}$ for
  An0n, by $P_\text{noEM}(1-P_\text{noEM})$ for Xn0n and by $P_\text{noEM}^2$
  for 0n0n.
\item \texttt{GAMMA\_AA\_FILE}, the probability $\Gamma_{AA}(b)$ that the two
  nuclei do not collide hadronically. It depends on the cross section
  $\sigma_{NN}$. The default file uses $\sigma_{NN} = 92$~mb , \texttt{input/WS\_photon\_flux/Gamma\_AA\_sigma90.85.dat} uses the
  $90.85$~mb of the Starlight tables. One make a new file with
  \texttt{src/make\_gamma\_aa.py}.
\end{itemize}
The Woods--Saxon parameters are $R = 6.49$~fm and $d = 0.54$~fm, the same as
in the Starlight tables.

\section{How the flux gets into the calculation}

In \texttt{src/photon\_flux.cpp} there are two functions that give $dN/d\omega$:

\begin{itemize}
\item Flux that depends on $b$ (\texttt{PL}, \texttt{WS}):
  \texttt{photon\_flux(b, qp, par)} gives
  $2\pi b\,\dfrac{d^3N}{d\omega\,d^2b}\,\Gamma_{AA}(b)\,P_\text{channel}(b)$.
\item Flux already integrated over $b$ (\texttt{EFF}, \texttt{STARLIGHT}):
  \texttt{effective\_photon\_flux(qp, par)} gives $\dfrac{dN}{d\omega}$.
\end{itemize}

\noindent The first one still has to be integrated over $b$; the code does
that as one of the variables of its Monte Carlo integration (VEGAS). For a
point-like nucleus the flux at fixed $b$ is
\begin{equation}
  \frac{d^3N}{d\omega\,d^2b} = \frac{\alpha Z^2}{\pi^2\,\omega}\,
  \frac{\eta^2}{b^2}\,K_1^2(\eta), \qquad \eta = z_\gamma\,m_N\,b ,
\end{equation}
and the Woods--Saxon and effective fluxes are built from the same quantity.
The Starlight table contains $\ln(dN/dz_\gamma)$, so the code takes the exponential
and multiplies by $2/\sqrt{s_{NN}}$, inside
\texttt{effective\_photon\_flux()}. 

The cross section is computed in \texttt{src/integrands.cpp}. There the code
asks which kind of flux it has (with \texttt{is\_effective\_flux()} in
\texttt{src/def.hpp}) and then uses the value in the same way for every
flux:
\begin{verbatim}
double flux = is_effective_flux(par) ? effective_photon_flux(qp, p)
                                     : photon_flux(b, qp, p);
...
double kernel = z * m*m * jac_qp * flux * g;      // for example, exclusive
\end{verbatim}
Everything else that depends on the photon energy ($z = p^+/q^+$, the
photon--nucleus energy $W^2$, $x_\mathbb{P}$, \dots) is computed here from
$q^+$, and it is the same for every flux. Then, the flux can be changed without modifying the integrands.

\subsection*{Careful: dipole impact parameter $b_d$}

The impact parameter $b_d$ is the position in the target nucleus
where the dipole (the charm quark--antiquark pair) interacts. The dipole
amplitude depends on $b_d$, so it is given as Glauber samples, one file per
value, \texttt{data/Pb/mve/glauber\_mve\_<}$b_d$\texttt{>} with
$b_d = 0, 2, \dots, 32$~GeV$^{-1}$ (17 files).

The tables (\texttt{EFF}, \texttt{STARLIGHT}) only remove the integral over the
flux $b$. The integral over $b_d$ is still needed for every flux:
\texttt{run\_scripts/run\_nucleus.sh} runs the D$^0$ program once for each
Glauber file, and \texttt{plotting\_scripts/D0.py} then computes
\begin{equation}
  \frac{d\sigma}{dy\,d^2p_{D^0\perp}} = 2\pi \int db_d\, b_d\;
  \frac{d\sigma(b_d)}{dy\,d^2p_{D^0\perp}}
\end{equation}
with Simpson's rule (\texttt{integrate\_over\_b}). So with a table flux, one D$^0$
run for a single Glauber file is not the cross section yet but one
point of this $b_d$ integral. 

\section{How to add another flux}\label{sec:own}

\subsection{A table in $z_\gamma$}

The easiest
way is to copy what is done for the Starlight table
(\texttt{init\_starlight\_flux()} and the \texttt{STARLIGHT} part of
\texttt{effective\_photon\_flux()} in \texttt{src/photon\_flux.cpp}):
\begin{enumerate}
\item Write a function \texttt{init\_<name>\_flux(channel)} that reads your
  table. In \texttt{effective\_photon\_flux()}, add a part for your flux that
  returns \textbf{$dN/d\omega$}. If your table has $dN/dz_\gamma$, multiply by
  $2/\sqrt{s_{NN}}$. If it has $z_\gamma\,dN/dz_\gamma$, divide by $z_\gamma$ first. If it has a
  logarithm, take the exponential first.
\item Add the name of your flux to \texttt{is\_effective\_flux()} in
  \texttt{src/def.hpp}. Otherwise the code thinks your flux still depends on
  $b$ and integrates it over $b$ a second time.
\item Call \texttt{init\_<name>\_flux()} where the other fluxes are set up:
  in the Pb+Pb part of \texttt{src/main.cpp} and \texttt{src/main\_xpom.cpp},
  and in \texttt{src/scan\_flux.cpp} to check it.
\end{enumerate}
No need to add factors of $z_\gamma$, $\omega$ or $1/q^+$ for the flux in the
integrands. 


\subsection{A flux that depends on $b$}

Add a part to \texttt{photon\_flux(b, qp, par)} that returns
$2\pi b\,(d^3N/d\omega\,d^2b)$. The factors $\Gamma_{AA}(b)$ and
$P_\text{channel}(b)$ are applied after it, as for \texttt{PL} and
\texttt{WS}. Do not add this flux to \texttt{is\_effective\_flux()}.

\section{Scan flux}\label{sec:checks}


 The program \texttt{build/bin/scan\_flux}
prints the flux as $dN/dz_\gamma$ for many values of $z_\gamma$, using exactly the same
functions as the D$^0$ calculation:
\begin{verbatim}
./build/bin/scan_flux STARLIGHT AnAn An0n     # or EFF, PL, WS, or your flux
\end{verbatim}
If the is a table, the output must give the table back. To compare several fluxes in one plot, run
\texttt{run\_scripts/run\_flux\_scan.sh} and then
\texttt{plotting\_scripts/flux\_comparison.py} with the same
$\sigma_{NN}$.




\end{document}
